\documentclass[12pt,a4paper]{report}
\usepackage[utf8]{inputenc}
\usepackage[french]{babel}
\usepackage[T1]{fontenc}
\usepackage{graphicx}
\usepackage{geometry}
\usepackage{hyperref}
\usepackage{booktabs}
\usepackage{fancyhdr}
\usepackage{titlesec}
\usepackage{float}
\usepackage{verbatim}

% Configuration des Marges
\geometry{top=2.5cm, bottom=2.5cm, left=2.5cm, right=2.5cm}

% En-têtes et Pieds de page
\pagestyle{fancy}
\fancyhf{}
\fancyhead[L]{StudyMate AI - Rapport de Projet}
\fancyhead[R]{Tuteur IA Pédagogique}
\fancyfoot[C]{\thepage}

\titleformat{\chapter}[hang]{\normalfont\large\bfseries}{\thechapter.}{1em}{}

\begin{document}

% ─── PAGE DE GARDE ───────────────────────────────────────────
\begin{titlepage}
    \centering
    \vspace*{1.5cm}
    {\huge\bfseries Université Technologique}\\[0.5cm]
    {\large Rapport de Projet de Fin d'Études}\\[1.5cm]
    
    \rule{\linewidth}{0.5mm}\\[0.4cm]
    {\Huge \bfseries StudyMate AI}\\[0.2cm]
    {\large \bfseries Tuteur Pédagogique Intelligent Fondé sur un Agent IA Autonome}\\[0.4cm]
    \rule{\linewidth}{0.5mm}\\[1.5cm]
    
    \noindent
    \begin{minipage}[t]{0.45\linewidth}
        \textbf{Rédigé par :}\\
        Étudiant en Intelligence Artificielle
    \end{minipage}
    \hfill
    \begin{minipage}[t]{0.45\linewidth}
        \textbf{Supervisé par :}\\
        Comité Pédagogique d'Évaluation
    \end{minipage}\\[3cm]
    
    {\large \today}
    \vfill
\end{titlepage}

% ─── TABLE DES MATIÈRES ──────────────────────────────────────
\tableofcontents
\newpage

% ─── RÉSUMÉ ───────────────────────────────────────────────────
\chapter*{Résumé}
\addcontentsline{toc}{chapter}{Résumé}
StudyMate AI est un assistant pédagogique destiné à faciliter la révision des étudiants. Il transforme des documents PDF en ressources de travail structurées : résumés, quiz, flashcards et planning de révision.

Le projet se concentre sur l'essentiel : rendre les cours plus accessibles et organiser l'apprentissage sans complexité technique.
\newpage

% ─── 1. INTRODUCTION ─────────────────────────────────────────
\chapter{Introduction}
StudyMate AI vise à aider les étudiants à mieux réviser en automatisant la transformation des supports de cours.

L'application propose une interface simple et des services concrets pour réduire le temps passé à préparer les révisions et augmenter l'efficacité de l'étude.

% ─── 2. PROBLÉMATIQUE ────────────────────────────────────────
\chapter{Problématique}
Les étudiants doivent souvent assimiler des documents volumineux sans outil adapté. Ils manquent de temps pour :
\begin{itemize}
    \item synthétiser le cours,
    \item s'auto-évaluer,
    \item planifier leurs sessions de révision.
\end{itemize}

Une application qui concentre ces fonctions peut améliorer la qualité des révisions.

% ─── 3. OBJECTIFS DU PROJET ──────────────────────────────────
\chapter{Objectifs du projet}
Les objectifs principaux de StudyMate AI sont :
\begin{itemize}
    \item extraire le contenu des PDF,
    \item générer des résumés pédagogiques,
    \item proposer des quiz et des flashcards,
    \item construire un planning de révision simple.
\end{itemize}

Le but est de fournir une aide pratique sans surcharge.

% ─── 4. ARCHITECTURE ET FONCTIONNEMENT ───────────────────────
\chapter{Architecture et fonctionnement}
Le système combine une interface web, un agent IA et des modules spécialisés.

\section{Interface utilisateur}
L'interface Gradio permet de charger un PDF, puis d'accéder aux fonctions de résumé, quiz, flashcards et planning.

\section{Agent IA}
L'agent pilote les actions et utilise des prompts pour guider la génération du contenu.

\section{Modules métiers}
Les modules du projet sont :
\begin{itemize}
    \item le lecteur PDF,
    \item le générateur de quiz,
    \item le générateur de flashcards,
    \item le planificateur de révision.
\end{itemize}

% ─── 5. TECHNOLOGIES UTILISÉES ───────────────────────────────
\chapter{Technologies utilisées}
Le projet repose sur :
\begin{itemize}
    \item Python,
    \item Gradio,
    \item LangChain,
    \item l'API Gemini,
    \item PyMuPDF.
\end{itemize}

Ces choix simplifient le développement tout en offrant des fonctionnalités robustes.

% ─── 6. DESCRIPTION DES OUTILS ───────────────────────────────
\chapter{Description des outils}
Chaque outil apporte une fonctionnalité clé.

\section{PDF Reader}
Il extrait le texte des fichiers PDF pour le rendre exploitable.

\section{Quiz Generator}
Il génère des questions à choix multiples liées au contenu du cours.

\section{Flashcard Generator}
Il crée des cartes de révision courtes pour mémoriser les notions essentielles.

\section{Planner Tool}
Il construit un planning adapté à la durée de révision disponible.

% ─── 7. INTERFACE UTILISATEUR ────────────────────────────────
\chapter{Interface utilisateur}
L'interface est épurée pour permettre un usage rapide. Les étudiants peuvent lancer les actions principales en quelques clics.

La navigation se fait par onglets et chaque résultat est présenté de manière lisible.

% ─── 8. RÉSULTATS ET OBSERVATIONS ────────────────────────────
\chapter{Résultats et observations}
StudyMate AI permet de produire des ressources utiles pour la révision :
\begin{itemize}
    \item résumés synthétiques,
    \item quiz de vérification,
    \item flashcards pour la mémorisation,
    \item planning simple.
\end{itemize}

Ces résultats montrent que l'outil répond aux besoins essentiels sans complexité inutile.

% ─── 9. LIMITES ET AMÉLIORATIONS ────────────────────────────
\chapter{Limites et améliorations}
Les limites principales sont :
\begin{itemize}
    \item dépendance à l'API externe,
    \item qualité variable selon le format du PDF,
    \item extraction difficile pour des documents mal structurés.
\end{itemize}

Des améliorations futures peuvent porter sur le stockage des sessions et la gestion des PDF volumineux.

% ─── 10. CONCLUSION ─────────────────────────────────────────
\chapter{Conclusion}
StudyMate AI offre une solution simple pour aider les étudiants à réviser. Il transforme des supports PDF en ressources structurées et facilite l'organisation de la révision.

La conception modulaire du projet permet d'ajouter facilement de nouvelles fonctions à l'avenir.

\end{document}